% Appendix C - Photochromic background
\section{Photochromic Materials: Background for the Patch}\label{app:photochromic}

This appendix explains, in simple terms, how the patch changes color and what that means
for Uvie. It is background for the prototype in Section~\ref{sec:prototype} and for the risks in
Section~\ref{sec:business-model}. It does not report new experiments.

\subsection{What Photochromism Is}
Photochromism is a reversible change in a material's light absorption, and therefore its color,
that is triggered by light. Several families of molecules behave this way, for example
spiropyrans, spirooxazines, azobenzenes, and diarylethenes. Spiropyrans are among the most
widely studied of these, and have become a molecule of choice for building responsive materials
\cite{klajn2014}.

\subsection{How a Spiropyran Changes Color}
A spiropyran has two forms (Figure~\ref{fig:photochromic}). In its closed form (SP) it is almost
colorless. Ultraviolet light, around 365\,nm in laboratory work, breaks one bond in the molecule and
opens the ring, which produces the colored merocyanine form (MC). The ring-opening follows
first-order kinetics. The reverse reaction happens again in the dark, driven by heat, and visible light
can speed it up \cite{klajn2014}. This is why a patch that turns dark in the sun slowly lightens
once the UV is gone.

\begin{figure}[ht]
  \centering
  \begin{tikzpicture}[
      st/.style={draw, rounded corners, minimum width=4.6cm, minimum height=1.5cm,
                 align=center, font=\small},
      >=Stealth]
    \node[st, fill=white] (sp) {\textbf{Closed form (SP)}\\almost colorless};
    \node[st, fill=violet!35, right=3.4cm of sp] (mc) {\textbf{Open form (MC)}\\colored};
    \draw[->, thick] ([yshift=0.25cm]sp.east) -- node[above, font=\small] {UV light} ([yshift=0.25cm]mc.west);
    \draw[->, thick] ([yshift=-0.25cm]mc.west) -- node[below, font=\small, align=center] {heat or\\visible light} ([yshift=-0.25cm]sp.east);
  \end{tikzpicture}
  \caption{Reversible color change of a spiropyran (simplified; after \cite{klajn2014}).}
  \label{fig:photochromic}
\end{figure}

\subsection{What This Means for Uvie}
The mechanism explains both the strengths and the weak points of the patch.
\begin{itemize}[itemsep=0pt]
  \item \textbf{It reacts to UV, not to sunscreen.} The color change is caused by UV reaching the
    molecule. A patch under sunscreen darkens more slowly only if the sunscreen blocks enough UV.
    In a laboratory test of ten commercial photochromic indicators, the devices changed color within
    about a minute at UV index 10 even under SPF 50+ sunscreen, and their color saturated at moderate
    UV levels \cite{aguilera2024}. A dark patch should therefore prompt the user to check and reapply,
    and it should not be read as a precise measure of remaining protection.
  \item \textbf{It fades again.} Because the colored form returns to the closed form on its own, the
    patch state depends on time and temperature as well as UV. How fast this happens in the warm
    Vietnamese climate has to be measured for our own formulation.
  \item \textbf{The host material matters.} In a solid patch the molecules are embedded in a matrix,
    which may influence speed and color contrast. The ink and the patch base are therefore part of
    the design, and their effect has to be tested.
\end{itemize}

\subsection{Characterization Plan}
Table~\ref{tab:photochromic-plan} lists what we need to measure before making any claim about the
patch. None of these measurements has been completed yet.

\begin{table}[ht]
\centering\small
\caption{Planned characterization of the photochromic patch}\label{tab:photochromic-plan}
\begin{tabular}{@{}L{3.6cm}L{6.4cm}L{4.1cm}@{}}
  \toprule
  \textbf{Property} & \textbf{How to measure it} & \textbf{Why it matters}\\
  \midrule
  Response time & Time to reach a set luminance drop under a UV lamp or measured sunlight & Sets how soon the app can warn\\
  \addlinespace
  Saturation & Luminance at increasing UV index values & Shows where the patch stops distinguishing levels\\
  \addlinespace
  Fade time & Time to return to the light state in shade and in warm conditions & Defines how long a reading stays valid\\
  \addlinespace
  Effect of sunscreen & Same test with no sunscreen and with several SPF levels & Tests the central claim of the product\\
  \addlinespace
  Repeatability & Several light--dark cycles on the same patch & Shows whether the patch ages\\
  \addlinespace
  Skin safety & Review of pigment, solvent, and adhesive by a laboratory & Needed before use on skin\\
  \bottomrule
\end{tabular}
\end{table}

For the phone measurement, each patch should be photographed with the same phone, distance, and
lighting, and the average luminance of the central area recorded. These values would replace the
first-estimate threshold of 120 used by the app (Figure~\ref{fig:confidence}).
